%% ===================================================================
\section{The cross section at order \texorpdfstring{$g^4$}{g4}}\label{sec:xsec}
%% ===================================================================

We now insert the coefficients into the amplitudes~\eqref{eq:amp1} and~\eqref{eq:amp2} and expand to order $g^4$. As announced after Eq.~\eqref{eq:LOrelab}, from here on the amplitude sits at $w$ and its complex conjugate at $w'$. Every term is accompanied by its complex conjugate, denoted ``$+$c.c.'', which is obtained by $w\leftrightarrow w'$, renaming the remaining primed and unprimed integration variables, reversing the order of all valence charges, and complex conjugating the numbers. We write $\mf A\equiv\mf A_2$ and $\mf A^\dagger\equiv\mf A^\dagger_2$, and a bar denotes rotation by the target, Sec.~\ref{sec:xsecdef}. Repeated colour and transverse-vector indices are summed.

Throughout, the colour operators of the individual terms are written with the adjoint index order that follows from Eq.~\eqref{eq:Srot}, $\bar a^{\dagger d}=U^{db}a^{\dagger b}$, wherever this matters for the derivations of Sec.~\ref{sec:LL}. In the term-by-term expressions of Sec.~\ref{sec:finite} the order of the two adjoint indices of $U$ is not tracked; there we use the shorthand $U^{bc}\equiv U^{cb}$. The leading-log results of Sec.~\ref{sec:LL} are derived for a general orthogonal $U$ and do not rely on this shorthand.

\subsection{Virtual terms: one gluon in the final state}\label{sec:groupI}

The one-gluon amplitude to order $g^3$ is the LO amplitude plus five classes of corrections:
\begin{equation}
\Big[U^{ab}(w)\,\mf A^{(1)b}_i(k^+,-w)-\bar{\mf A}^{(1)a}_i(k^+,-w)\Big]+\sum_{X={\rm I}}^{\rm V}\mathcal V^a_{X,i}(w),
\label{eq:amp1NLO}
\end{equation}
with
\begin{align}
\mathcal V^a_{{\rm I},i}(w)={}&-\bar{\mf A}^{(1)a}_i(k^+,-w)\,\mf N_2+U^{ab}(w)\,\bar{\mf N}_2\,\mf A^{(1)b}_i(k^+,-w),\label{eq:VI}\\
\mathcal V^a_{{\rm II},i}(w)={}&U^{ab}(w)\,\mf A^{(3)b}_i(k^+,-w)-\bar{\mf A}^{(3)a}_i(k^+,-w)-\big[\bar{\mf N}_2,\bar{\mf A}^{(1)a}_i(k^+,-w)\big],\label{eq:VII}\\
\mathcal V^a_{{\rm III},i}(w)={}&\int\frac{dp^+}{2\pi}\int_x\bar{\mf B}^{\dagger ba}_{3ji}(k^+,w;p^+,-x)\Big[-\bar{\mf A}^{(1)b}_j(p^+,-x)+U^{bc}(x)\,\mf A^{(1)c}_j(p^+,-x)\Big],\label{eq:VIII}\\
\mathcal V^a_{{\rm IV},i}(w)={}&2\int\frac{dp^+}{2\pi}\int_x\bar{\mf A}^{\dagger(1)b}_j(p^+,-x)\Big[U^{ac}(w)U^{bd}(x)\,\mf B^{cd}_{2ji}(k^+,-w;p^+,-x)-\bar{\mf B}^{ba}_{2ji}(k^+,-w;p^+,-x)\Big],\label{eq:VIV}\\
\mathcal V^a_{{\rm V},i}(w)={}&2\int\frac{dp^+}{2\pi}\int_{x,z}\Big[-\bar{\mf C}^{\dagger bca}_{1jki}(k^+,w-z;p^+,z-x)\,\bar{\mf B}^{bc}_{2jk}(k^+-p^+,-z;p^+,-x)\notag\\
&\hspace{6em}+\bar{\mf C}^{\dagger dca}_{1jki}(k^+,w-z;p^+,z-x)\,U^{db}(x)U^{ce}(z)\,\mf B^{be}_{2jk}(k^+-p^+,-z;p^+,-x)\Big].\label{eq:VV}
\end{align}
In $\mathcal V_{\rm II}$ we used the unitarity relation~\eqref{eq:unitA3X} for $\bar{\mf A}^{(3)\dagger}_1$ and moved its $\mf B$ and $\mf C$ terms into $\mathcal V_{\rm III}$--$\mathcal V_{\rm V}$. The physical content of the five classes is:
\begin{itemize}
\item[I:] the normalization of the incoming ($\mf N_2$) and outgoing ($\bar{\mf N}_2$) dressed states, multiplying the emission after and before the target, respectively;
\item[II:] the one-loop correction to the emission vertex: gluon self-energy, vertex corrections and the two- and three-charge terms of $\mf A^{(3)}$;
\item[III:] a gluon of momentum $p^+$ emitted before ($U\mf A$) or after ($\bar{\mf A}$) the target is converted into the measured gluon by the final-state dressing $\bar\Omega^\dagger$, through its interaction with the rotated valence charges ($\bar{\mf B}_3^\dagger$);
\item[IV:] two gluons produced before the target ($\mf B_2$), one of which is absorbed by a valence charge after the target ($\bar{\mf A}^\dagger$);
\item[V:] two gluons produced before the target recombine into the measured gluon after it ($\bar{\mf C}^\dagger$).
\end{itemize}
The Group~I contribution to the cross section is the product with the conjugate LO amplitude,
\begin{equation}
\frac{d^3N^{\rm (I)}}{d^3k}=\frac{1}{(2\pi)^3}\int_{w,w'}e^{-ik\cdot(w'-w)}\Big[-\bar{\mf A}^{\dagger(1)a}_i(k^+,-w')+\mf A^{\dagger(1)b'}_i(k^+,-w')\,U^{ab'}(w')\Big]\sum_{X={\rm I}}^{\rm V}\mathcal V^a_{X,i}(w)+{\rm c.c.}
\label{eq:CS1}
\end{equation}
We call the five terms Rows G1-I to G1-V.

\paragraph*{The diagonal term G1-I$'$.} $\mf B_3$ is defined with the window $|p^+-k^+|<\Lambda$ excluded, since there the outgoing gluon is indistinguishable from the incoming one, see Sec.~\ref{sec:singleglue}. In the product $\bar\Omega^\dagger\Omega$, however, the measured gluon produced before the target can emit and reabsorb a soft gluon \emph{after} the target. In the coefficient language this is the diagonal part $p^+\to k^+$ of $\bar{\mf B}^\dagger_3$ acting on the LO amplitude. We denote it by Row G1-I$'$. Its leading-log form is derived from the leading-log wave function in Sec.~\ref{sec:LL} and listed in App.~\ref{app:LLrows}. It supplies the JIMWLK terms in which both derivatives act on the Wilson line of the measured gluon.
\note{The finite part of G1-I$'$ requires the diagonal element of $\bar{\mf B}_3^\dagger$ beyond leading log, which is not yet computed.}
%TODO: compute Row G1-I' (diagonal region of B3-bar-dagger) beyond LL.

\subsection{Real terms: two gluons in the final state}\label{sec:groupII}

The two-gluon amplitude~\eqref{eq:amp2} at order $g^2$, for the measured gluon $(a,i,k^+)$ at $w$ and the unobserved gluon $(b,j,p^+)$ at $z$, is $\frac{2}{\sqrt{2\pi}}$ times
\begin{equation}
\mathfrak B^{ab}_{ij}(w,z)+\mathfrak X^{ab}_{1,ij}(w,z)+\mathfrak X^{ab}_{2,ij}(w,z)+\mathfrak C^{ab}_{ij}(w,z),
\label{eq:amp2NLO}
\end{equation}
with
\begin{align}
\mathfrak B^{ab}_{ij}={}&U^{ac}(w)\,\mf B^{bc}_{2ji}(k^+,-w;p^+,-z)-U^{\dagger bc}(z)\,\bar{\mf B}^{ca}_{2ji}(k^+,-w;p^+,-z),\label{eq:frakB}\\
\mathfrak X^{ab}_{1,ij}={}&\frac12\Big[U^{\dagger bc}(z)\,\bar{\mf A}^{(1)c}_j(p^+,-z)\,\bar{\mf A}^{(1)a}_i(k^+,-w)-U^{ac}(w)\,\bar{\mf A}^{(1)d}_j(p^+,-z)\,U^{\dagger db}(z)\,\mf A^{(1)c}_i(k^+,-w)\Big],\label{eq:frakX1}\\
\mathfrak X^{ab}_{2,ij}={}&\frac12\Big[U^{\dagger bc}(z)\,\bar{\mf A}^{(1)a}_i(k^+,-w)\,\bar{\mf A}^{(1)c}_j(p^+,-z)-\bar{\mf A}^{(1)a}_i(k^+,-w)\,\mf A^{(1)b}_j(p^+,-z)\Big],\label{eq:frakX2}\\
\mathfrak C^{ab}_{ij}={}&-\int_y\Big[\bar{\mf C}^{dac}_{1jik}(p^++k^+,y-w;p^+,w-z)\,U^{\dagger bd}(z)U^{ce}(y)\,\mf A^{(1)e}_k(p^++k^+,-y)\notag\\
&\hspace{4em}-U^{\dagger bd}(z)\,\bar{\mf C}^{dac}_{1jik}(p^++k^+,y-w;p^+,w-z)\,\bar{\mf A}^{(1)c}_k(p^++k^+,-y)\Big].\label{eq:frakC}
\end{align}
$\mathfrak B$ contains the two-gluon component of the projectile produced before ($U\mf B_2$) or after ($\bar{\mf B}_2$) the target. $\mathfrak X_{1,2}$ are the two operator orderings of two independent emissions, which follow from the unitarity relation~\eqref{eq:unitB1X} for $\mf B_1^\dagger$. $\mathfrak C$ is the splitting of a single gluon into two after the target (fragmentation of the scattered gluon); the parent has momentum $p^++k^+$.

The squared amplitude splits into diagonal products (Group~II) and cross products (Group~III):
\begin{align}
\frac{d^3N^{\rm (II)}}{d^3k}={}&\frac{4}{(2\pi)^3}\int\frac{dp^+}{2\pi}\int_{z,w,w'}e^{-ik\cdot(w'-w)}\Big\{\mathfrak B^\dagger\mathfrak B+\mathfrak X_1^\dagger\mathfrak X_1+\mathfrak X_2^\dagger\mathfrak X_2+\mathfrak C^\dagger\mathfrak C\Big\},\label{eq:CS2}\\
\frac{d^3N^{\rm (III)}}{d^3k}={}&\frac{4}{(2\pi)^3}\int\frac{dp^+}{2\pi}\int_{z,w,w'}e^{-ik\cdot(w'-w)}\Big\{\mathfrak X_1^\dagger\mathfrak B+\mathfrak X_2^\dagger\mathfrak B+\mathfrak C^\dagger\mathfrak B+\mathfrak X_2^\dagger\mathfrak X_1+\mathfrak C^\dagger\mathfrak X_1+\mathfrak C^\dagger\mathfrak X_2\Big\}+{\rm c.c.}\label{eq:CS3}
\end{align}
Here $\mathfrak Y^\dagger\mathfrak Z\equiv\sum_{a,b,i,j}[\mathfrak Y^{ab}_{ij}(w',z)]^\dagger\,\mathfrak Z^{ab}_{ij}(w,z)$. We call the four terms of Eq.~\eqref{eq:CS2} Rows G2-I to G2-IV and the six terms of Eq.~\eqref{eq:CS3} Rows G3-I to G3-VI, in the order written.

\subsection{Overview of the terms}\label{sec:rows}

Table~\ref{tab:rows} summarizes the content of every row. Two properties are important for what follows.

\paragraph*{Longitudinal integrals.} After the transverse structure is fixed, each row is a single integral over the plus momentum $p^+$ of the unobserved (real) or virtual gluon. Large logarithms arise only where $p^+$ is strongly ordered with respect to $k^+$:
\begin{itemize}
\item for $\Lambda<p^+\ll k^+$ (region T) the integrand behaves as $\mathfrak T/p^+$ and gives $\log(k^+/\Lambda)$;
\item for $k^+\ll p^+<\vee$ (region P) it behaves as $\mathfrak P/p^+$ and gives $\log(\vee/k^+)$.
\end{itemize}
Configurations with $p^+\sim k^+$ give no large logarithm. The total rapidity interval is
\begin{equation}
\delta Y\equiv\log\frac{\vee}{\Lambda}=\delta Y_T+\delta Y_P,\qquad \delta Y_T\equiv\log\frac{k^+}{\Lambda},\qquad \delta Y_P\equiv\log\frac{\vee}{k^+}.
\label{eq:dYdef}
\end{equation}
When a logarithm appears with an argument involving $\vee-k^+$ or $k^+\pm\Lambda$, we always rewrite it with the identities
\begin{equation}
\log\frac{\Lambda}{k^+}=-\log\frac{\vee}{\Lambda}+\log\frac{\vee}{k^+},\qquad \log\frac{\vee-k^+}{\Lambda}=\log\frac{\vee}{\Lambda}+\log\Big(1-\frac{k^+}{\vee}\Big),
\label{eq:logid}
\end{equation}
and keep $\log(1-k^+/\vee)$ exactly, since $\vee$ is finite.

\paragraph*{Number of charges.} The rows contain products of two, three or four valence charges. Only the two-charge terms survive at leading log (Sec.~\ref{sec:cancel}); the finite parts of Sec.~\ref{sec:finite} contain three-charge terms as well.

\begin{table}[!tbp]
\centering
\renewcommand{\arraystretch}{1.25}
\begin{tabular}{llL{5.7cm}cc}
\toprule
Row & coefficients & process & charges & $p^+$ range\\
\midrule
G1-I & $\mf N_2\mf A$ & normalization of the dressed state & 4 & $(\Lambda,\vee)$\\
G1-II & $\mf A^{(3)}$, $[\bar{\mf N}_2,\bar{\mf A}]$ & one-loop emission vertex (self-energy, $\beta_0$) & 2, 3, 4 & $(\Lambda,\vee)$\\
G1-III & $\bar{\mf B}_3^\dagger\mf A$ & gluon $p^+$ converted into the measured gluon after the target & 3 & $(\Lambda,\vee)$, $|p^+-k^+|>\Lambda$\\
G1-IV & $\bar{\mf A}^\dagger\mf B_2$ & two gluons before the target, one absorbed after it & 3 & $(\Lambda,\vee-k^+)$\\
G1-V & $\bar{\mf C}^\dagger\mf B_2$ & two gluons before the target recombine after it (loop across the target) & 2, 3 & $(0,k^+)$\\
G1-I$'$ & diagonal $\bar{\mf B}_3^\dagger$ & measured gluon emits and reabsorbs a soft gluon after the target & 2 & $|p^+-k^+|<\Lambda$\\
\midrule
G2-I & $\mathfrak B^\dagger\mathfrak B$ & two gluons from the projectile wave function, before or after the target; initial-state splitting & 2, 3, 4 & $(\Lambda,\vee-k^+)$\\
G2-II & $\mathfrak X_1^\dagger\mathfrak X_1$ & two independent emissions & 4 & $(\Lambda,\vee-k^+)$\\
G2-III & $\mathfrak X_2^\dagger\mathfrak X_2$ & two independent emissions & 4 & $(\Lambda,\vee-k^+)$\\
G2-IV & $\mathfrak C^\dagger\mathfrak C$ & splitting of the scattered gluon; fragmentation & 2 & $(\Lambda,\vee-k^+)$\\
\midrule
G3-I, II & $\mathfrak X_{1,2}^\dagger\mathfrak B$ & interference & 3 & $(\Lambda,\vee-k^+)$\\
G3-III & $\mathfrak C^\dagger\mathfrak B$ & interference of initial- and final-state splitting & 2, 3 & $(\Lambda,\vee-k^+)$\\
G3-IV & $\mathfrak X_2^\dagger\mathfrak X_1$ & interference of the two orderings & 4 & $(\Lambda,\vee-k^+)$\\
G3-V, VI & $\mathfrak C^\dagger\mathfrak X_{1,2}$ & interference & 3 & $(\Lambda,\vee-k^+)$\\
\bottomrule
\end{tabular}
\caption{The rows of the NLO cross section, Eqs.~\eqref{eq:CS1}, \eqref{eq:CS2} and~\eqref{eq:CS3}. ``Charges'' is the number of valence charges $\rho$ in the products, before reordering. The $p^+$ range is that of the unobserved or virtual gluon.}
\label{tab:rows}
\end{table}
